\documentclass[leqno, a4paper, 12pt]{jsarticle}
\usepackage{amssymb}
\usepackage{amsthm}
\usepackage{amsmath}
\usepackage{comment}
\usepackage{url}

\usepackage{enumitem}
%\usepackage{showkeys}


%定理環境 thmはあまり使わずpropをなるべく使う。
%主張は基本的に証明の中で使い、そのため番号を付けない。
%注意環境を付けてみた。
\newtheorem{thm}{定理}
\newtheorem{prop}[thm]{命題}
\newtheorem{cor}[thm]{系}
\newtheorem{lem}[thm]{補題}
\newtheorem{df}[thm]{定義}
\newtheorem{exam}[thm]{例}
\newtheorem{fact}[thm]{事実}
\newtheorem*{claim}{主張}
\newtheorem*{cau}{注意}
\newtheorem*{rmk}{注意}
\renewcommand{\proofname}{証明}
%矢印たち。
%\toは\rightarrowと同じ。\getsは\leftarrowと同じ。同値は\iff
\newcommand{\To}{\Longrightarrow}
\newcommand{\Gets}{\LongLeftarrow}
%黒板太字
\newcommand{\nn}{\mathbb{N}}
\newcommand{\zz}{\mathbb{Z}}
\newcommand{\qq}{\mathbb{Q}}
\newcommand{\rr}{\mathbb{R}}
\newcommand{\cc}{\mathbb{C}}
%無限大
\newcommand{\mgn}{\infty}
%差集合は\setminusで統一する。等号の否定は\neq
%真部分集合は\subsetneqq　偏微分は\partial
%ディスプレイスタイル。\dfracでディスプレイ分数
\newcommand{\dpst}{\displaystyle}

\title{論文メモ
\large{\\--- 不動点定理，主にCaristiの ---}
}
\author{ナンブ　キトラ}
\date{\today}
\begin{document}
\maketitle

本棚を整理したいので，
溜まっている論文のメモを書いて未来への手紙とすることにする．

以下の論文の束は
Caristiの不動点定理を調べて印刷したものだと思う．
Caristiの不動点定理というのはBanachの不動点定理の拡張で以下のようなものである．
\begin{thm}
$(X, d)$を完備距離空間とする．
$T\colon X\to X$を単に写像とし，
$f\colon X\to [0, \infty)$
は下半連続関数とする．
もしも任意の$x\in X$について
\[
d(x, T(x))\le f(x)-f(T(x))
\]
が成り立つのならば，
$T$は不動点をもつ．
つまり$T(x)=x$となる$x\in X$
が存在する．
\end{thm}

論文\cite{MR0394329}, \cite{MR0399968}
はいわゆる原論文である．


論文
\cite{MR0229101}
はブラウワーとか角谷とかの不動点定理の話で
上の不動点定理とは無関係だと思う．

論文
\cite{MR0237724} 
と
\cite{MR0236331}は
$[0, 1]$からそれ自身への二つの写像
$f$と$g$が可換つまり$f\circ g=g\circ f$
であっても共通の不動点を持つとは限らない例を構成している．


論文
\cite{MR0458359}
は下半連続関数を使って
距離空間の完備性を
特徴付けしている．

文献\cite{1}, 
\cite{2}, 
\cite{3}
は日本語で書かれた不動点定理の解説などである．



\section*{参考文献たち}

\renewcommand{\refname}{日本語の文献}
	
\begin{thebibliography}{99}
\bibitem[和1]{1}
鈴木智成, 
「CARISTIの不動点定理の拡張定理の別証明」, 
数理解析研究所講究録, 
No. 1415 (2002), 146--151, 
http://hdl.handle.net/2433/26259


\bibitem[和2]{2}
鈴木智成, 
「距離空間における不動点定理の4つ の分類」
数理解析研究所講究録, 
No. 1753(2011),  52--57
http://hdl.handle.net/2433/171179

\bibitem[和3]{3}
木島洋一
「距離空間における不動点定理」, 
横浜経営研究, 
巻 18, 号 2, p. 99--102, 発行日 1997-09-01, 
横浜国立大学経営学会, 
http://hdl.handle.net/10131/758
\end{thebibliography}

\renewcommand{\refname}{英文参考文献}
%READ!
%myplain is the same to amsplain style. 
\nocite{*}
\bibliographystyle{myplaindoidoi}
\bibliography{fixed.bib}



\end{document}